\documentclass[11pt,a4paper]{report}

% ==================== PACKAGES ====================
\usepackage[margin=2.5cm]{geometry}
\usepackage[utf8]{inputenc}
\usepackage[T1]{fontenc}
\usepackage{graphicx}
\usepackage{booktabs}
\usepackage{longtable}
\usepackage{amsmath}
\usepackage{amssymb}
\usepackage{float}
\usepackage{caption}
\usepackage{subcaption}
\usepackage{xcolor}
\usepackage{listings}
\usepackage{fancyhdr}
\usepackage[hidelinks]{hyperref}
\usepackage[backend=biber,style=numeric,sorting=none]{biblatex}
\addbibresource{references.bib}

% ==================== PAGE STYLE ====================
\pagestyle{fancy}
\fancyhf{}
\fancyhead[L]{Laporan Teknis -- MagangHub Batch 2 2026}
\fancyhead[R]{\thepage}
\renewcommand{\headrulewidth}{0.4pt}

% ==================== LISTINGS STYLE (untuk kode) ====================
\lstset{
    basicstyle=\ttfamily\small,
    breaklines=true,
    frame=single,
    numbers=left,
    numberstyle=\tiny,
    columns=fullflexible,
    backgroundcolor=\color{gray!5}
}

% ==================== JUDUL ====================
\title{
    \Large \textbf{Laporan Teknis Analisis Data dan Dashboard} \\
    \large MagangHub Batch 2 2026
} 
\author{Nama Penulis / Tim \\ Institusi / Program Studi}
\date{\today}

% ====================================================
\begin{document}

\maketitle

\begin{abstract}
Tuliskan ringkasan 150--250 kata yang mencakup: (1) latar belakang singkat proyek analisis data lowongan magang MagangHub Batch 2 2026, (2) tujuan analisis, (3) metode yang digunakan (EDA deskriptif, competition ratio, analisis geografis, dsb.), dan (4) ringkasan temuan utama serta luaran berupa dashboard.
\end{abstract}

\tableofcontents
\listoffigures
\listoftables

% ==================== BAB 1 ====================
\chapter{Pendahuluan}

\section{Latar Belakang}
Jelaskan konteks MagangHub, program magang nasional/bersertifikat, dan mengapa analisis data terhadap lowongan magang ini penting dilakukan.

\section{Rumusan Masalah (Business Questions)}
Rujuk daftar lengkap pada \texttt{business\_questions.md} atau salin ringkasannya di sini, misalnya sebagai daftar bernomor BQ1--BQ9 (dan BQ10--BQ14 bila digunakan).

\section{Tujuan}
Jelaskan tujuan analisis, misalnya memahami karakteristik lowongan, tingkat persaingan, distribusi geografis, serta menyusun dashboard.

\section{Ruang Lingkup dan Batasan}
Jelaskan batasan data (periode Batch 2 2026), batasan interpretasi (competition ratio bukan probabilitas diterima, dsb.), dan batasan metode (analisis deskriptif, bukan inferensial/kausal).

% ==================== BAB 2 ====================
\chapter{Data}

\section{Sumber Data}
Jelaskan sumber data (scraping/ekspor dari platform MagangHub), periode pengambilan data, dan jumlah data mentah.

\section{Deskripsi Variabel}
Tabel deskripsi kolom dataset.

\begin{table}[H]
\centering
\caption{Deskripsi Variabel Dataset}
\begin{tabular}{ll}
\toprule
\textbf{Kolom} & \textbf{Deskripsi} \\
\midrule
posisi & Nama posisi/lowongan magang \\
instansi & Nama instansi penyelenggara \\
jenis\_instansi & Pemerintah / Swasta \\
jurusan\_1..n & Daftar jurusan yang dapat mendaftar \\
lokasi\_kabupaten & Kabupaten/kota penempatan \\
lokasi\_provinsi & Provinsi penempatan \\
jenjang & Jenjang pendidikan yang dipersyaratkan \\
hari\_kerja & Jumlah hari kerja per minggu \\
kuota & Jumlah kuota peserta \\
pelamar & Jumlah pelamar tercatat \\
label\_peluang & Kategori peluang dari platform \\
\bottomrule
\end{tabular}
\end{table}

\section{Proses Data Cleaning dan Preparation}
Jelaskan tahapan parsing HTML/teks menjadi tabular, penanganan missing value, normalisasi nama wilayah, dan transformasi kolom jurusan ke format long.

% ==================== BAB 3 ====================
\chapter{Metodologi}

\section{Pendekatan Analisis}
Jelaskan pendekatan yang dipakai: statistik deskriptif, perhitungan competition ratio, analisis geografis, analisis korelasi/scatter plot, dsb.

\section{Tools dan Teknologi}
Sebutkan tools (Python/pandas, visualisasi dengan matplotlib/plotly, dashboard dengan Streamlit/Power BI/Looker Studio, dsb.).

\section{Alur Kerja Analisis}
Sertakan diagram alur (pipeline) dari data mentah hingga dashboard.

\begin{figure}[H]
\centering
% \includegraphics[width=0.8\textwidth]{figures/pipeline.png}
\caption{Alur kerja analisis data MagangHub Batch 2 2026}
\label{fig:pipeline}
\end{figure}

% ==================== BAB 4 ====================
\chapter{Hasil dan Pembahasan}

\section{Karakteristik Umum Lowongan (BQ1)}
Sajikan hasil overview: total lowongan, total kuota, total pelamar, rata-rata, dan distribusinya.

\section{Tingkat Persaingan (BQ2)}
Sajikan hasil competition ratio secara umum dan distribusinya.

\section{Perbandingan Instansi Pemerintah vs Swasta (BQ3)}

\section{Distribusi Geografis: Lowongan dan Pelamar (BQ4)}

\section{Distribusi Geografis: Tingkat Persaingan (BQ5)}

\section{Analisis Posisi (BQ6)}

\section{Analisis Jurusan (BQ7)}

\section{Distribusi Jenjang Pendidikan (BQ8)}

\section{Hubungan Kuota dan Pelamar (BQ9)}

\section{Analisis Tambahan (BQ10--BQ14)}
Opsional -- isi bila pertanyaan tambahan digunakan (pola jadwal kerja, validasi label peluang, top instansi, korelasi jurusan/jenjang terhadap persaingan).

% ==================== BAB 5 ====================
\chapter{Rancangan Dashboard}

\section{Struktur dan Fitur Dashboard}
Jelaskan struktur dashboard: ringkasan umum, peta distribusi wilayah, perbandingan instansi, top posisi/jurusan, distribusi jenjang, scatter plot kuota vs pelamar.

\section{Visualisasi Utama}
Sertakan contoh tangkapan layar/wireframe dashboard.

\begin{figure}[H]
\centering
% \includegraphics[width=0.9\textwidth]{figures/dashboard_overview.png}
\caption{Contoh tampilan dashboard MagangHub Batch 2 2026}
\label{fig:dashboard}
\end{figure}

\section{Implementasi}
Jelaskan tools implementasi dashboard dan cara mengaksesnya.

% ==================== BAB 6 ====================
\chapter{Kesimpulan dan Rekomendasi}

\section{Kesimpulan}
Ringkas temuan utama dari setiap business question.

\section{Rekomendasi}
Berikan rekomendasi bagi pemangku kepentingan (mis. calon pelamar, penyelenggara program, pembuat kebijakan).

\section{Keterbatasan}
Jelaskan keterbatasan data dan metode analisis.

% ==================== DAFTAR PUSTAKA ====================
\printbibliography[title={Daftar Pustaka}]

% ==================== LAMPIRAN ====================
\appendix
\chapter{Lampiran}

\section{Tabel Tambahan}

\section{Cuplikan Kode}
\begin{lstlisting}[language=Python, caption={Contoh perhitungan competition ratio}]
df['competition_ratio'] = df['pelamar'] / df['kuota']
\end{lstlisting}

\end{document}